\documentclass[11pt,a4paper]{article}
\usepackage[margin=1in]{geometry}
\usepackage[spanish]{babel}
\usepackage[utf8]{inputenc}
\usepackage[T1]{fontenc}
\usepackage{array}
\usepackage{booktabs}
\usepackage{enumitem}
\usepackage{hyperref}
\usepackage{siunitx}
\usepackage{tabularx}
\setlength{\parskip}{6pt}
\setlength{\parindent}{0pt}
\sisetup{output-decimal-marker={,}}

\title{Analisis tecnico: multi-rep por sesion y limpieza de UI}
\author{MOOVIA Sensor Lab}
\date{\today}

\begin{document}
\maketitle

\section*{Resumen ejecutivo}
Esta iteracion amplifica el analisis posterior al \emph{Stream Off}: el sistema deja de asumir una unica repeticion implicita y pasa a construir un bloque especifico de \texttt{repAnalysis} sobre el tramo activo ya segmentado. El objetivo no es hacer conteo en vivo, sino ofrecer un resumen util y trazable por sesion con:
\begin{itemize}[leftmargin=14pt]
  \item numero de repeticiones completas;
  \item lista de reps cerradas con inicio, apex, fin y metricas propias;
  \item posible rep parcial final, reportada aparte y no contada;
  \item limpieza de textos visibles en la app movil para eliminar mojibake y dejar labels legibles.
\end{itemize}

\section{Motivacion}
Hasta ahora la app y el web-test ya podian describir un tramo activo unico: velocidad media propulsiva, altura, lateral, salud de captura y resumen del endpoint activo/asentado. Eso era suficiente para sesiones de una sola repeticion, pero no para series reales.

El cambio de enfoque es importante:
\begin{itemize}[leftmargin=14pt]
  \item la metrica global de sesion sigue existiendo y no se rompe;
  \item por encima de ella se anade un analisis por repeticion, calculado \textbf{solo offline}, una vez terminada la captura;
  \item el sistema deja claro que una serie puede contener reps completas y una ultima rep incompleta.
\end{itemize}

\section{Estado previo}
Antes de este cambio ya existia una maquina de estados interna con transiciones del tipo \texttt{IDLE -> MOVING -> APEX -> RETURN}, pero ese flujo servia como apoyo para el seguimiento del gesto, no como feature de producto. En concreto, faltaban cuatro piezas:
\begin{enumerate}[leftmargin=16pt]
  \item un contrato publico de \texttt{repAnalysis};
  \item una lista de reps con indices y tiempos;
  \item metricas por rep desacopladas de las metricas globales de sesion;
  \item una UI que mostrara el conteo y limpiara los textos rotos en la app movil.
\end{enumerate}

\section{Heuristica elegida}
La heuristica implementada trabaja sobre el \textbf{tramo activo ya segmentado}, no sobre toda la sesion. Esto evita mezclar colocacion inicial, reposo y cola quieta con la dinamica real del levantamiento.

\subsection*{Regla de ciclo completo}
La repeticion se define como \textbf{ciclo completo}:
\begin{itemize}[leftmargin=14pt]
  \item si la serie arranca subiendo, la rep se interpreta como \texttt{valle -> apex -> vuelta al valle};
  \item si la serie arranca bajando, la rep se interpreta como \texttt{pico -> valle -> vuelta al pico}.
\end{itemize}

\subsection*{Señal vertical usada}
La deteccion usa el eje \(Z\) ya reconstruido por el pipeline compartido, junto con una version suavizada de \(v_z\). No se exige estado estacionario duro para separar reps, de modo que dos reps seguidas sin pausa larga tambien puedan contarse.

\subsection*{Parametros}
Los parametros fijados en la implementacion son:
\begin{itemize}[leftmargin=14pt]
  \item ventana de suavizado de \SI{35}{ms} sobre \(z\) y \(v_z\);
  \item excursion minima de \SI{0,08}{m};
  \item tolerancia de retorno de \SI{0,04}{m};
  \item duracion minima de \SI{250}{ms};
  \item duracion maxima de \SI{8000}{ms};
  \item confirmacion del cambio de signo de velocidad durante al menos \SI{60}{ms}.
\end{itemize}

\section{Modelo de datos nuevo}
La salida compartida \texttt{SessionAnalysisSummary} incorpora un bloque \texttt{repAnalysis}. Sobre ese bloque se publican dos tipos nuevos:
\begin{itemize}[leftmargin=14pt]
  \item \texttt{RepetitionSummary}: describe una rep concreta con indices, tiempos, direccion, estado de completitud, confianza y metricas;
  \item \texttt{RepAnalysisSummary}: resume toda la serie con \texttt{repCount}, lista de reps, \texttt{partialRep}, VMP media de serie y mejor rep.
\end{itemize}

Las metricas por rep son:
\begin{itemize}[leftmargin=14pt]
  \item \texttt{meanPropulsiveVelocity};
  \item \texttt{peakVerticalVelocity};
  \item \texttt{peakLinearAcc};
  \item \texttt{maxHeight};
  \item \texttt{netHeight};
  \item \texttt{maxLateral};
  \item \texttt{finalLateral}.
\end{itemize}

\section{Diferencia frente a las metricas globales}
La capa multi-rep no sustituye a las metricas globales del tramo activo; las complementa.

\begin{table}[h!]
\centering
\small
\begin{tabularx}{\textwidth}{>{\raggedright\arraybackslash}p{3.2cm}X X}
\toprule
Bloque & Que mide & Para que sirve \\
\midrule
\texttt{movementMetrics} & resumen del tramo activo completo & ver la sesion como un unico gesto estabilizado \\
\texttt{repAnalysis.reps[i]} & metrica de una rep concreta & comparar reps dentro de una misma serie \\
\texttt{partialRep} & ultima rep no cerrada & avisar de que hubo gesto sin retorno suficiente \\
\bottomrule
\end{tabularx}
\end{table}

Esto evita romper compatibilidad con la app, el web-test y los JSON existentes, al tiempo que habilita una lectura mas deportiva y mas util para series reales.

\section{Cambios visibles en UI}
\subsection*{Web-test}
El web-test ahora muestra:
\begin{itemize}[leftmargin=14pt]
  \item tarjeta \texttt{Repetition Analysis};
  \item tabla con \#, direccion, duracion, VMP, velocidad pico, aceleracion pico, altura maxima, lateral maxima y confianza;
  \item marcadores de \texttt{start}, \texttt{apex} y \texttt{end} sobre la grafica \(Z(t)\).
\end{itemize}

\subsection*{App movil}
La app movil ahora muestra, tras \emph{Detener stream}:
\begin{itemize}[leftmargin=14pt]
  \item reps completas;
  \item VMP media de la serie;
  \item mejor rep;
  \item una card por rep con duracion, VMP, velocidad pico, aceleracion pico, altura maxima y lateral maxima;
  \item rep incompleta si existe.
\end{itemize}

Ademas se limpiaron los labels visibles que estaban afectados por mojibake. Botones y mensajes quedan en texto simple, por ejemplo:
\begin{itemize}[leftmargin=14pt]
  \item \texttt{Iniciar stream};
  \item \texttt{Detener stream};
  \item \texttt{Reset IMU};
  \item \texttt{Exportar JSON RAW};
  \item \texttt{Conectado} / \texttt{Desconectado}.
\end{itemize}

\section{Validacion con los JSON V3}
Aunque los V3 siguen siendo sesiones de gesto practicamente unico, son utiles para validar el nuevo flujo:
\begin{itemize}[leftmargin=14pt]
  \item \texttt{V3-moovia-session-1774637899200.json}: la heuristica detecta \textbf{1 rep completa} con direccion \texttt{up-first}. La forma \emph{sube un poco antes de bajar} sigue visible y la rep queda resumida con VMP, altura y lateral propias.
  \item \texttt{V3-moovia-session-1774637821525.json}: la heuristica no fuerza un falso positivo; registra \textbf{0 reps completas} y una \texttt{partialRep down-first}. Esto es coherente con una sesion que no vuelve lo suficiente al nivel inicial de la rep dentro de la tolerancia configurada.
\end{itemize}

Ese comportamiento es importante: el contador no inventa reps cerradas cuando la trayectoria IMU-only no ofrece evidencia suficiente de retorno.

\section{Limitaciones}
Este cambio mejora mucho la lectura de serie, pero mantiene los limites fisicos del sistema actual:
\begin{itemize}[leftmargin=14pt]
  \item sigue siendo un algoritmo \textbf{offline/postprocesado}, no conteo en vivo;
  \item depende de la calidad de la reconstruccion vertical del tramo activo;
  \item una deriva vertical grande puede convertir una rep real en \texttt{partialRep};
  \item la lateral sigue siendo relativa e IMU-only;
  \item el sistema no se debe vender como contador perfecto de reps multi-eje en cualquier ejercicio.
\end{itemize}

En otras palabras: el nuevo bloque multi-rep es muy util para sesiones verticales estructuradas, pero sigue heredando el techo de observabilidad IMU-only descrito en la documentacion V3 previa.

\section{Casos esperados de uso}
La heuristica esta pensada para soportar estos escenarios:
\begin{itemize}[leftmargin=14pt]
  \item \textbf{1 rep aislada}: debe devolver \texttt{repCount = 1} si la rep cierra bien;
  \item \textbf{3 reps seguidas sin pausa larga}: deben seguir contandose porque no se exige \texttt{isStatState} entre reps;
  \item \textbf{3 reps con pausa breve}: deben detectarse como reps independientes;
  \item \textbf{ultima rep incompleta}: debe aparecer en \texttt{partialRep} y no sumarse al contador.
\end{itemize}

\section*{Conclusion}
Esta iteracion hace dos cosas valiosas a la vez:
\begin{enumerate}[leftmargin=16pt]
  \item convierte el analisis compartido en una base real para series con multiples reps;
  \item limpia la UI movil para que el resultado sea legible y usable por cualquier persona del equipo.
\end{enumerate}

No resuelve el limite IMU-only, pero si organiza mejor la informacion y evita que el sistema trate toda la sesion como una sola repeticion implita. Ese paso era necesario para que la siguiente fase del proyecto pueda comparar reps, detectar parciales y acercarse a una herramienta de entrenamiento realmente util.

\end{document}
